% !TeX root = ../../../main.tex
\chapter{Geometric Class Field Theory}\label{section:gcft}

The ramification theorem of Chapter~2 is the key input for the proof of
\Cref{thm:GCFT_reduced}. Retain the notation and hypotheses of that theorem:
$C$ is a smooth, projective, geometrically connected curve over the perfect
field $k$, $\mdls$ is a modulus, $U=C\setminus\mdls$, $\Lambda$ is a finite
commutative ring whose cardinality is invertible in $k$, and $\cF$ is a
rank-one $\Lambda$-local system on $U$ whose ramification is bounded by
$\mdls$. Our goal is to descend $\cF^{[d]}$ along
\[
    \Phi_d\colon U^{(d)}\longrightarrow\PicCm[d]
\]
for all sufficiently large $d$.

Set $G=\Lambda^\times$ and let
$\cP=\operatorname{Fr}(\cF)$ be the frame torsor of $\cF$. By
\Cref{prop:torsor_module_equivalence}, together with its compatibility with
pullback, symmetric powers, and ramification, it is enough to construct a
$G$-torsor $\cQ_d$ on $\PicCm[d]$ such that
\[
    \Phi_d^*\cQ_d\cong\cP^{[d]}.
\]
We therefore work with $G$-torsors throughout the proof and return to
rank-one local systems only in the final step.

When $\mdls=0$, the projective-space fibers of $\Phi_d$ give the descent
directly. When $\mdls>0$, $\Phi_d$ is not proper and its fibers are affine,
so we instead use Takeuchi's extension to a projective space bundle
\[
    \widetilde\Phi_d\colon
    \Cmod{d}
    \longrightarrow \PicCm[d],
\]
where $\Cmod{d}$ is an open subscheme of the blowup of
$C^{(d)}$ along the locus obtained by adding the modulus.

The single-point calculation of Chapter~2, combined with compatibility
under products and addition of divisors, first shows that
$\cP^{[d]}$ is
tamely ramified along the boundary
\[
    H=\Cmod{d}\setminus U^{(d)}.
\]
Passing to the geometric generic fiber, we then show that $\cP^{[d]}$ extends
uniquely across $H$; the extension then descends along the projective-space
bundle $\widetilde\Phi_d$ to $\PicCm[d]$.

The remainder of the chapter makes these steps precise.
\Cref{section:generalized-picard-and-abel-jacobi} recalls the rigidified
Picard scheme and the Abel--Jacobi morphism.
\Cref{section:BlowupOfSymmetricPowerOfCurves} introduces Takeuchi's
compactification.
\Cref{section:tame-ramification-full-boundary} proves tameness along its
boundary, and \Cref{section:extension-across-boundary} extends the torsor
across that boundary. Finally, \Cref{section:gcft-descent} completes the
descent and returns from torsors to rank-one local systems.

% !TeX root = ../../../main.tex
\section{The Rigidified Picard Scheme and the Abel--Jacobi Map}\label{section:generalized-picard-and-abel-jacobi}

Our presentation of the rigidified Picard scheme and the Abel--Jacobi
morphism follows Guignard \cite[\S4]{Guignard2018} and
Takeuchi \cite[\S3]{takeuchi2018blow}. For background on the classical
theory of Jacobians and relative Picard schemes, we refer respectively to
Milne \cite{milneAV} and Bosch--L\"utkebohmert--Raynaud
\cite[Chapter~9]{BLR1990}. Precise references for the individual results
are given where they are used.

\subsection*{The rigidified Picard scheme}
Although the main theorem is over the fixed perfect field $k$, the
rigidified Picard construction is most naturally stated over an arbitrary
base scheme $S$.
Let $S$ be a scheme and let $f:C\to S$ be a projective smooth morphism whose geometric fibers are connected and of dimension 1. Let $\mdls$ be a modulus on $C$, defined as an effective Cartier divisor of $C/S$
(i.e., a closed subscheme of $C$ which is finite flat of finite presentation over $S$).
We denote the projection $C \times_S T \to T$ by $\operatorname{pr}$ for any $S$-scheme $T$.

Let $d$ be an integer. For an $S$-scheme $T$, we consider the set of data $(\cL, \psi)$ where:
\begin{itemize}
    \item $\cL$ is an invertible sheaf of degree $d$ on $C_T$.
    \item $\psi: \Oo_{\mdls_T} \xrightarrow{\sim} \cL|_{\mdls_T}$ is a trivialization of $\cL$ along the modulus.
\end{itemize}
Two such pairs $(\cL, \psi)$ and $(\cL', \psi')$ are said to be isomorphic if there exists an isomorphism of invertible sheaves $\alpha\colon \cL \to \cL'$ such that the following diagram commutes:
\[
\begin{tikzcd}
& \Oo_{\mdls_T} \arrow[dl, "\psi"'] \arrow[dr, "\psi'"] & \\
\cL|_{\mdls_T} \arrow[rr, "\alpha|_{\mdls_T}"'] & & \cL'|_{\mdls_T}
\end{tikzcd}
\]
We define the presheaf $\operatorname{Pic}_{C,\mdls}^{d,\mathrm{pre}}$ on
$\operatorname{Sch}/S$ by assigning to $T$ the set of isomorphism classes
of such pairs. Let $\PicCm[d]$ denote the \'etale sheafification of this
presheaf.

We use the following structural facts, recorded in
\cite[\S3, fundamental properties (1)--(4)]{takeuchi2018blow}. When $\mdls$ is everywhere
positive, representability and smoothness of the full rigidified Picard
presheaf are proved in \cite[Proposition~4.8]{Guignard2018}; when
$\mdls=0$, the corresponding results for the ordinary relative Picard
scheme are given in \cite[\S9.3, Theorem~1]{BLR1990}.
\begin{enumerate}
    \item $\PicCm[d]$ is represented by an $S$-scheme. If $\mdls$ is
    faithfully flat over $S$, the presheaf
    $\operatorname{Pic}_{C,\mdls}^{d,\mathrm{pre}}$ is already an \'etale
    sheaf.
    \item $\PicCm[0]$ is a smooth commutative group $S$-scheme with
    geometrically connected fibers.
    \item For any $d$, $\PicCm[d]$ is a $\PicCm[0]$-torsor.
\end{enumerate}
We call $\PicCm[0]$ the \emph{generalized Jacobian} of $C$ with modulus
$\mdls$. When $\mdls=0$, it is the ordinary Jacobian, and
$\PicCm[d]=\PicC[d]$.

The rigidified Picard schemes are naturally related as the modulus varies.
Indeed, let $\mdls_1\subset\mdls_2$. Restriction of rigidifications defines
a natural morphism of \'etale sheaves
\[
\rho_{\mdls_2,\mdls_1}\colon
\operatorname{Pic}_{C,\mdls_2}^d
\longrightarrow
\operatorname{Pic}_{C,\mdls_1}^d.
\]
When $\mdls_2$ is finite over $S$, this morphism is an epimorphism
\cite[\S3]{takeuchi2018blow}.
Taking $\mdls_1=0$ gives, for every
modulus $\mdls$, a natural epimorphism of \'etale sheaves
\[
\rho_{\mdls,0}\colon \PicCm[d]\longrightarrow\PicC[d].
\]

\subsection*{The Abel--Jacobi Morphism and Its Fibers}\label{section:AbelJacobi}


Let $U = C \setminus \mdls$ be the complement of the modulus in $C$.
The effective Cartier divisors of degree $d$ which are prime to $\mdls$ are parameterized by the symmetric power $\Sym_S^d (U) = U^{(d)}$ over $S$ (see \cite[Proposition~4.12]{Guignard2018} and \cite[Chapter~III, Theorem~3.13]{milneAV}).
For any such divisor $D \in U^{(d)}$, the associated line bundle $\Oo_C(D)$ admits a canonical trivialization along $\mdls$.
Specifically, the canonical section $1_D$ is regular and non-vanishing on $\mdls$ because $\operatorname{supp}(D) \cap \operatorname{supp}(\mdls) = \emptyset$.
Its restriction is a nowhere-vanishing section on $\mdls$, thereby
determining a trivialization
$\psi_D:\Oo_{\mdls}\xrightarrow{\sim}\Oo_C(D)|_{\mdls}$.
This is functorial in families and yields a morphism from the symmetric power
to the rigidified Picard scheme over $S$:
\begin{equation}
    \Phi_d\colon U^{(d)} \to \PicCm[d], \quad D \mapsto [(\Oo_C(D), \psi_D)],
\end{equation}

The results of Guignard and Takeuchi used below require the degree to be
sufficiently large.\footnote{The cohomological reason for this bound is that,
for a degree-$d$ line bundle $\cL$, it ensures that both $\cL$ and
$\cL(-\mdls)$ have fiberwise degree at least $2g-1$. Their first cohomology
then vanishes by Serre duality and a degree argument; Riemann--Roch together
with cohomology and base change implies that the corresponding pushforwards
are locally free of the expected ranks and commute with base change. See
\cite[Proposition~4.3]{Guignard2018}. The additional inequality
$d\geq\deg\mdls$ ensures that $d-\deg\mdls$ is nonnegative, which is relevant
when $g=0$.} Let $g$ denote the genus of $C$, viewed as a locally
constant function on $S$. From this point on, we impose the standing
condition
\begin{equation} \label{equation:degree-condition}
    d \geq \max\{2g - 1 + \deg \mdls, \deg \mdls\}.
\end{equation}
This is Takeuchi's condition \cite[\S3, (3.1)]{takeuchi2018blow} and is, in
particular, strong enough for the application of
\cite[Proposition~4.14]{Guignard2018} below. If $S$ is quasi-compact, the
locally constant functions $g$ and $\deg\mdls$ take only finitely many
values; hence a single integer $d$ satisfying this condition exists.

Suppose that $d$ satisfies \Cref{equation:degree-condition}. When $\mdls=0$
and $C\to S$ admits a section, the Abel--Jacobi morphism
\[
    C^{(d)}\longrightarrow\PicC[d]
\]
is a projective space bundle. In particular, it is proper and smooth with
geometrically connected fibers.

When $\mdls>0$, Guignard \cite[Proposition~4.14]{Guignard2018} proves the
stronger statement that
\[
    \Phi_d\colon U^{(d)}\longrightarrow\PicCm[d]
\]
is, for the fppf topology, a torsor under the additive group scheme of a
vector bundle of rank $d-\deg\mdls-g+1$ over $\PicCm[d]$. In particular,
$\Phi_d$ is fppf-locally isomorphic over $\PicCm[d]$ to an affine space of
relative dimension $d-\deg\mdls-g+1$; it is smooth and surjective, with
geometrically connected fibers.

We now return to the standing situation $S=\Spec(k)$, which remains in
force for the rest of the chapter. From Guignard's result, we retain only
the following fiberwise consequence, stated together with the classical
unramified case.
\begin{theorem}\label{theorem:abel-jacobi-fibers}
Assume that $S=\Spec(k)$ and that $d$ satisfies
\Cref{equation:degree-condition}. Let
$\bar s=\Spec\Omega\to\PicCm[d]$ be a geometric point, where $\Omega$ is
separably closed. Then the geometric fiber of the Abel--Jacobi morphism
\[
    \Phi_d\colon U^{(d)}\longrightarrow\PicCm[d]
\]
at $\bar s$ is isomorphic to
\[
\begin{cases}
\mathbb{A}_{\Omega}^{d - \deg \mdls - g + 1} & \text{if } \mdls > 0, \\
\mathbb{P}_{\Omega}^{d - g} & \text{if } \mdls = 0.
\end{cases}
\]
Thus every geometric fiber is an affine space when $\mdls>0$ and a
projective space when $\mdls=0$.
\end{theorem}
\begin{proof}
For $\mdls>0$, this follows from
\cite[Proposition~4.14]{Guignard2018}. For $\mdls=0$, see
\cite[Proposition~2.1.4]{Toth2011}.
Both cases are also stated in
\cite[Propositions~3.13--3.14]{tendler2015geometricclassfieldtheory}.
\end{proof}



% !TeX root = ../../../main.tex
\section{Takeuchi's Compactification of the Abel--Jacobi Morphism}\label{section:BlowupOfSymmetricPowerOfCurves}
We recall that our objective is to descend the $G$-torsor $\cP^{[d]}$ from
$U^{(d)}$ to $\PicCm[d]$ along the Abel--Jacobi map $\Phi_d$:

\[
\begin{tikzcd}
 \cP^{[d]} \arrow[d, purple]\\
 U^{(d)} \arrow[r, "\Phi_d"] & \PicCm[d]
\end{tikzcd}
\]


(Here, the purple arrow emphasizes that $\cP^{[d]}$ is a torsor on the
\'etale site of $U^{(d)}$.)

In the ramified case $\mdls>0$, \Cref{theorem:abel-jacobi-fibers}
identifies the geometric fibers of $\Phi_d$ with affine spaces. Since
$\Phi_d$ is not proper, the proper-smooth fundamental-group argument
available in the unramified case cannot be applied directly. Takeuchi
\cite[\S3]{takeuchi2018blow} overcomes this difficulty by extending $\Phi_d$
to a projective-space bundle. We now recall the ambient blowup and the
properties of the resulting compactification that will be used below.

Let $\mdls = \sum_{i=1}^s n_iP_i$, where the $P_i$ are distinct closed
points of $C$, be a modulus on $C$, and let $d$ satisfy
\Cref{equation:degree-condition}.
Takeuchi \cite[\S3, before (3.6)]{takeuchi2018blow} defines $Z_0=Z_0(\mdls,d)$ as the closed subscheme of
$C^{(d)}$ given by the image of the closed immersion
$C^{(d-\deg\mdls)}\to C^{(d)}$ that adds $\mdls$.
He also defines $X_{\mdls,d}$ as the blowup of $C^{(d)}$ along $Z_0$.
Let $E_0=E_{\mdls,d}=Z_0(\mdls,d)\times_{C^{(d)}}X_{\mdls,d}$ be the
exceptional divisor. It is irreducible of codimension~$1$, and we let
$\eta_0=\eta_{\mdls,d}$ be its generic point.


Diagrammatically:
\[\begin{tikzcd}
\overline{\{\eta_0 \}} =  E_0 \arrow[r] \arrow[d] & X_{\mdls, d} \arrow[d, "\pi"] \\
Z_0 \arrow[r, hook]           & C^{(d)}
\end{tikzcd}\]

Incorporating $U^{(d)}$, the torsor $\cP^{[d]}$, and the Abel--Jacobi map,
we have:
\[
\begin{tikzcd}
    & & \cP^{[d]} \arrow[d, purple]\\
\overline{\{\eta_0 \}} =  E_0 \arrow[r] \arrow[d] &
X_{\mdls, d} \arrow[d, "\pi"]  & \arrow[l, hook'] \arrow[dl, hook'] U^{(d)} \arrow[r, "\Phi_d"] & \PicCm[d] \\
Z_0 \arrow[r, hook]           & C^{(d)}
\end{tikzcd}
\]

We will not need Takeuchi's explicit fiber-product definition of the
compactification. The properties required below are summarized in the
following existence and structure statement, which follows from
\cite[\S3, (3.9), Proposition~3.7, and Corollary~3.8]{takeuchi2018blow}.

\begin{theorem}[Takeuchi]\label{theorem:takeuchi-compactification-theorem}
    Let $d$ satisfy \Cref{equation:degree-condition}. There exists an open
    subscheme $\Cmod{d}\subset X_{\mdls,d}$ containing $U^{(d)}$ and a
    morphism
    \[
        \widetilde\Phi_d\colon \Cmod{d}\longrightarrow\PicCm[d]
    \]
    extending $\Phi_d\colon U^{(d)}\to\PicCm[d]$, such that $\Cmod{d}$ is a
    projective space bundle over $\PicCm[d]$.
    Furthermore, the complement of $U^{(d)}$ in $\Cmod{d}$ is the scheme-theoretic intersection
    \[
        E_0 \times_{X_{\mdls,d}} \Cmod{d}
        \cong Z_0 \times_{C^{(d)}} \Cmod{d}.
    \]
\end{theorem}



Diagrammatically we have:
\[
\begin{tikzcd}
    &  & \cP^{[d]} \arrow[d, purple]\\
    E_0 \times_{X_{\mdls,d}} \Cmod{d} \arrow[d] \arrow[r]  & \Cmod{d} \arrow[d, hook] \arrow[rd, "\widetilde\Phi_d"]  & U^{(d)} \arrow[l, hook'] \arrow[d, "\Phi_d"] \\
\overline{\{\eta_0 \}} =  E_0 \arrow[r] \arrow[d] &
X_{\mdls, d} \arrow[d, "\pi"]  &  \PicCm[d] &  \\
Z_0 \arrow[r, hook]           & C^{(d)}
\end{tikzcd}
\]

Takeuchi's description of the compactification identifies its boundary as a
relative hyperplane; see
\cite[Lemmas~3.2 and~3.3(1), Theorem~3.4, (3.9), and
Corollary~3.8]{takeuchi2018blow}. This gives the following fiberwise
description.

\begin{cor}[The boundary is a relative hyperplane]
\label{lemma:geometric-generic-fiber-takeuchi}
The boundary
\[
    H=\Cmod{d}\setminus U^{(d)}
\]
is a smooth relative hyperplane divisor in the projective-space bundle
$\widetilde\Phi_d:\Cmod{d}\to\PicCm[d]$. Consequently, if
$\bar s=\Spec\Omega\to\PicCm[d]$ is any geometric point and
$r=d-\deg\mdls-g+1$, then there is an isomorphism of pairs\footnote{For pairs
$(X,D)$ and $(Y,E)$, an isomorphism of pairs is an isomorphism
$f\colon X\xrightarrow{\sim}Y$ such that
$D=X\times_Y E$ as closed subschemes of $X$. No canonical choice of $f$ is
asserted here.}
\[
    \left((\Cmod{d})_{\bar s},H_{\bar s}\right)
    \cong
    \left(\mathbb P^r_{\Omega},\mathbb P^{r-1}_{\Omega}\right),
\]
and hence $U^{(d)}_{\bar s}\cong\mathbb A^r_{\Omega}$.
\end{cor}

\begin{proof}
The formation of a projective bundle and of its relative hyperplane commutes
with arbitrary base change. Hence, after base change to $\bar s$, the
projective-space bundle $\Cmod{d}$ becomes $\mathbb P^r_{\Omega}$ and $H$
becomes a hyperplane $\mathbb P^{r-1}_{\Omega}$. Its complement is therefore
$\mathbb A^r_{\Omega}$.
\end{proof}



% !TeX root = ../../../main.tex
\section{Tame Ramification Along the Full Boundary}\label{section:tame-ramification-full-boundary}

The ramification calculation of Chapter~2 proves the required tameness when
the modulus is supported at a single point and the degree is the degree of
that modulus. We now remove these two restrictions. First, we pass from the
point-supported components of $\mdls$ to the full modulus in degree
$\deg\mdls$. We then compare the blowup at the point
$\mdls\in C^{(\deg\mdls)}$ with the blowup in degree $d$ along the locus of
divisors containing $\mdls$. This proves the following result.

\begin{theorem}\label{theorem:SymmetricPowerOfSheafIsTamelyRamified}
Let $G$ be a finite abelian group, and let $\cP$ be a $G$-torsor on $U$
whose ramification is bounded by $\mdls$. For every sufficiently large
integer $d$, the torsor $\cP^{[d]}$ is tamely ramified along
\[
H=\Cmod{d}\setminus U^{(d)}
  =E_0\times_{X_{\mdls,d}}\Cmod{d}
  \cong Z_0\times_{C^{(d)}}\Cmod{d}.
\]
\end{theorem}

To prove the theorem, we retain the blowup notation of Chapter~2. For every
nonzero effective submodulus $\ndls\leq\mdls$, we identify $\ndls$ with the
corresponding $k$-point of $C^{(\deg\ndls)}$, and put
\[
    X_{\ndls}=\Bl_{\ndls}\bigl(C^{(\deg\ndls)}\bigr),
    \qquad
    E_{\ndls}=X_{\ndls}\times_{C^{(\deg\ndls)}}\{\ndls\}.
\]
We write $\eta_{\ndls}$ for the generic point of $E_{\ndls}$. For the full
modulus in degree $d$, the center is
\[
Z_0=Z_0(\mdls,d)=\mdls+C^{(d-\deg\mdls)}\subset C^{(d)},
\]
where the right-hand side denotes the image of the closed immersion
$D\mapsto\mdls+D$.
The corresponding blowup, exceptional divisor, and generic point are denoted
by $X_{\mdls,d}$, $E_0=E_{\mdls,d}$, and
$\eta_0=\eta_{\mdls,d}$, as in
\Cref{section:BlowupOfSymmetricPowerOfCurves}.

\noindent\textbf{From point-supported moduli to the full modulus.}

We first show that if the relevant torsors are tame for two coprime
submoduli, then the corresponding torsor is tame for their sum. Iterating
this step over the point-supported components of $\mdls$ yields the required
statement for the full modulus in degree $\deg\mdls$.

To establish this implication, we use the following general comparison. Let
$X$ and $Y$ be smooth integral schemes over a field $k$, and let
$x\in X(k)$ and $y\in Y(k)$.
We denote the blowups of these schemes at the given points by
$\pi_X: \Bl_x(X) \to X$ and $\pi_Y: \Bl_y(Y) \to Y$. Furthermore, let $\pi_{X \times Y}: \Bl_{(x,y)}(X \times_k Y) \to X \times_k Y$
be the blowup of the product scheme at the point $(x,y)$.
We denote by $E_X, E_Y$, and $E_{X \times Y}$ the respective exceptional divisors, and let $\eta_X, \eta_Y$, and $\eta_{X \times Y}$ be their generic points.

Tame ramification is preserved under the contracted product of torsors after
pullback to a product, as follows.

\begin{prop}\label{prop:ramification-external-product}
Let $G$ be a finite abelian group. Suppose $\cP_X$ and $\cP_Y$ are
$G$-torsors defined on dense open subsets
$U_X\subset\Bl_x(X)$ and $U_Y\subset\Bl_y(Y)$ disjoint from the respective
exceptional divisors. There exist a dense open subset
\[
    \widetilde U\subset\Bl_{(x,y)}(X\times_kY)
\]
containing $\eta_{X\times Y}$ and morphisms
\[
    r_X:\widetilde U\longrightarrow\Bl_x(X),
    \qquad
    r_Y:\widetilde U\longrightarrow\Bl_y(Y),
\]
which send $\eta_{X\times Y}$ to $\eta_X$ and $\eta_Y$, respectively. On
the dense open subset
\[
    W=r_X^{-1}(U_X)\cap r_Y^{-1}(U_Y)\subset\widetilde U,
\]
define
\[
    \cP_{X\times Y}
      :=r_X^{-1}\cP_X
        \otimes^G
        r_Y^{-1}\cP_Y.
\]
If $\cP_X$ is tamely ramified at $\eta_X$ and $\cP_Y$ is tamely ramified
at $\eta_Y$, then $\cP_{X\times Y}$ is tamely ramified at
$\eta_{X\times Y}$.
\end{prop}

\begin{proof}
We first construct $r_X$ and $r_Y$; the torsors enter only after these
morphisms have been defined on the original blowups. Shrinking $X$ and $Y$
around $x$ and $y$ does not change the exceptional divisors or their local
rings at the generic points. The \`etale local description of smooth
morphisms \stackstag{054L} therefore gives \`etale morphisms
\[
a:X\longrightarrow\mathbb A_k^n,
\qquad
b:Y\longrightarrow\mathbb A_k^m
\]
carrying $x$ and $y$ to the origins; after shrinking further, we may assume
that $x$ and $y$ are the scheme-theoretic inverse images of the respective
origins. We write $x_1,\ldots,x_n$ and $y_1,\ldots,y_m$ for the corresponding
standard coordinates and also for their pullbacks to $X$ and $Y$.

Write
\[
    B_n=\operatorname{Bl}_0(\mathbb A_k^n),\qquad
    B_m=\operatorname{Bl}_0(\mathbb A_k^m),\qquad
    B_{n,m}=\operatorname{Bl}_{(0,0)}
        (\mathbb A_k^n\times_k\mathbb A_k^m),
\]
and denote their exceptional divisors by $E_n$, $E_m$, and $E_{n,m}$.
The blowup $B_{n,m}$ is the subscheme of
$(\mathbb A_k^n\times_k\mathbb A_k^m)\times_k\mathbb P_k^{n+m-1}$ defined by

\[\begin{cases}
x_i w_j = x_j w_i & 1 \le i, j \le n \\
y_k w_{n+l} = y_l w_{n+k} & 1 \le k, l \le m \\
x_i w_{n+j} = y_j w_i & 1 \le i \le n, \,\, 1 \le j \le m
\end{cases}\]
where $[w_1:\dots:w_{n+m}]$ are the homogeneous coordinates of
$\mathbb P_k^{n+m-1}$. The exceptional divisor $E_{n,m}$ is
isomorphic to $\mathbb P_k^{n+m-1}$.

Both $B_n\times_kB_m$ and $B_{n,m}$ are birational to
$\mathbb A_k^n\times_k\mathbb A_k^m$; their common dense open is
\[
    (\mathbb A_k^n\setminus\{0\})
      \times_k(\mathbb A_k^m\setminus\{0\}).
\]
Let
\[
    q_{n,m}\colon B_{n,m}
      \longrightarrow\mathbb P_k^{n+m-1}
\]
be the second projection. Define
\[
    \Omega_n=\bigcup_{i=1}^{n}D_+(w_i),
    \qquad
    \Omega_m=\bigcup_{j=1}^{m}D_+(w_{n+j}),
\]
and
\[
    \widetilde U_{n,m}
      :=q_{n,m}^{-1}(\Omega_n\cap\Omega_m)
      =\bigcup_{\substack{1\leq i\leq n\\1\leq j\leq m}}
        q_{n,m}^{-1}\!\left(D_+(w_iw_{n+j})\right).
\]
There is a morphism
\[
\begin{aligned}
f_{n,m}\colon\widetilde U_{n,m}&\longrightarrow
   B_n\times_kB_m,\\
((u,v),[w_1:\cdots:w_{n+m}])&\longmapsto
\bigl((u,[w_1:\cdots:w_n]),
      (v,[w_{n+1}:\cdots:w_{n+m}])\bigr).
\end{aligned}
\]
Thus we have a diagram
\[
\begin{tikzcd}
& \widetilde U_{n,m}
    \arrow[dl,"f_{n,m}"'] \arrow[dr,hook] & \\
B_n\times_kB_m && B_{n,m}.
\end{tikzcd}
\]
The restriction of $q_{n,m}$ induces an isomorphism
$E_{n,m}\xrightarrow{\sim}\mathbb P_k^{n+m-1}$. Since neither
$V_+(w_1,\ldots,w_n)$ nor $V_+(w_{n+1},\ldots,w_{n+m})$ contains the generic
point of $\mathbb P_k^{n+m-1}$, we have
$\eta_{n,m}\in\widetilde U_{n,m}$, where $\eta_{n,m}$ is the generic point
of $E_{n,m}$.
On the exceptional divisors, the restriction of $f_{n,m}$ is the
coordinate-block morphism
\[
\begin{aligned}
E_{n,m}\cap\widetilde U_{n,m}&\longrightarrow E_n\times_kE_m,\\
[w_1:\cdots:w_{n+m}]&\longmapsto
\bigl([w_1:\cdots:w_n],[w_{n+1}:\cdots:w_{n+m}]\bigr).
\end{aligned}
\]
This morphism is surjective and hence dominant. It therefore sends
$\eta_{n,m}$ to the generic point $(\eta_n,\eta_m)$ of $E_n\times_kE_m$.

We now return to the original schemes. By
\Cref{theorem:blowup_flat_base_change}, there are canonical identifications
\[
\begin{aligned}
\operatorname{Bl}_x(X)&\cong X\times_{\mathbb A_k^n} B_n,\\
\operatorname{Bl}_y(Y)&\cong Y\times_{\mathbb A_k^m} B_m,\\
B:=\operatorname{Bl}_{(x,y)}(X\times_kY)
  &\cong (X\times_kY)
      \times_{\mathbb A_k^n\times_k\mathbb A_k^m}B_{n,m}.
\end{aligned}
\]
Set
\[
    \widetilde U
      :=(X\times_kY)
        \times_{\mathbb A_k^n\times_k\mathbb A_k^m}\widetilde U_{n,m}.
\]
It is a dense open neighborhood of $\eta_{X\times Y}$ in $B$. Base change of
$f_{n,m}$ gives a morphism
\[
    f:\widetilde U\longrightarrow
      \operatorname{Bl}_x(X)\times_k\operatorname{Bl}_y(Y)
\]
that sends $\eta_{X\times Y}$ to $(\eta_X,\eta_Y)$. If $p_X$ and $p_Y$ are
the two projections from its target, set
\[
    r_X=p_X\circ f,
    \qquad
    r_Y=p_Y\circ f.
\]
These are the morphisms asserted in the statement. They fit into the
commutative diagram
\[
\begin{tikzcd}[column sep=large,row sep=large]
& \operatorname{Bl}_{(x,y)}(X\times_kY) & \\
& \widetilde U
    \arrow[u,hook] \arrow[dl,"r_X"'] \arrow[d,"f"] \arrow[dr,"r_Y"] & \\
\operatorname{Bl}_x(X)
& \operatorname{Bl}_x(X)\times_k\operatorname{Bl}_y(Y)
    \arrow[l,"p_X"] \arrow[r,"p_Y"']
& \operatorname{Bl}_y(Y).
\end{tikzcd}
\]
Over the common dense
open $(X\setminus\{x\})\times_k(Y\setminus\{y\})$, the morphism $f$ is the
evident identification; in particular, it is dominant. Hence
\[
    W=r_X^{-1}(U_X)\cap r_Y^{-1}(U_Y)
      =f^{-1}(U_X\times_k U_Y)
\]
is a nonempty open subset of the integral scheme $\widetilde U$, and is
therefore dense in $\widetilde U$ and in $B$.

It remains to compare the maps induced by $r_X$ and $r_Y$ on the exceptional
valuations. Let
\[
S=\mathcal O_{B,\eta_{X\times Y}},\qquad
R_X=\mathcal O_{\operatorname{Bl}_x(X),\eta_X},\qquad
R_Y=\mathcal O_{\operatorname{Bl}_y(Y),\eta_Y}.
\]
The morphisms $r_X$ and $r_Y$ induce local homomorphisms
\[
    R_X\longrightarrow S,
    \qquad
    R_Y\longrightarrow S.
\]
By base change, the exceptional divisors are identified with
\[
    E_X\cong\mathbb P_k^{n-1},\qquad
    E_Y\cong\mathbb P_k^{m-1},\qquad
    E_{X\times Y}\cong\mathbb P_k^{n+m-1},
\]
and the restrictions of $r_X$ and $r_Y$ are the two coordinate projections
\[
[w_1:\cdots:w_{n+m}]\longmapsto[w_1:\cdots:w_n],
\qquad
[w_1:\cdots:w_{n+m}]\longmapsto[w_{n+1}:\cdots:w_{n+m}].
\]
At the generic point of $E_{X\times Y}$ both coordinate blocks are nonzero.
On the chart $D_+(w_1w_{n+1})$, the functions $x_1$ and $y_1$ are
uniformizers in $R_X$ and $R_Y$, respectively, and each is also a
uniformizer in $S$. Hence both extensions have ramification index one. Their
residue-field extensions are
\[
k\left(\frac{x_2}{x_1},\ldots,\frac{x_n}{x_1}\right)
\subset
k\left(\frac{x_2}{x_1},\ldots,\frac{x_n}{x_1},
        \frac{y_1}{x_1},\ldots,\frac{y_m}{x_1}\right)
\]
and
\[
k\left(\frac{y_2}{y_1},\ldots,\frac{y_m}{y_1}\right)
\subset
k\left(\frac{x_1}{y_1},\ldots,\frac{x_n}{y_1},
        \frac{y_2}{y_1},\ldots,\frac{y_m}{y_1}\right),
\]
respectively. In particular, both are purely transcendental and hence
separable.

We now apply this comparison to the torsors. The extensions
$R_X\to S$ and $R_Y\to S$ are weakly unramified and have
separable residue-field extensions. Hence
\Cref{lemma:weakly_unramified_base_change} shows that the pullbacks of
$\cP_X$ and $\cP_Y$ to $\operatorname{Frac}(S)$ are tamely ramified.
By the definition above, their contracted product is the generic fiber of
$\cP_{X\times Y}$. Therefore
\Cref{lemma:bounding_ramification_of_contracted_product} shows that
$\cP_{X\times Y}$ is tamely ramified at $\eta_{X\times Y}$.
\end{proof}

\begin{cor}\label{cor:moduli-reduction}
Let $G$ be a finite abelian group, and let $\cP$ be a $G$-torsor on $U$.
Let $\ndls_1,\ndls_2\subset\mdls$ be two coprime nonzero submoduli.
Assume that $\cP^{[\deg\ndls_1]}$ and $\cP^{[\deg\ndls_2]}$ are tamely
ramified at $\eta_{\ndls_1}$ and $\eta_{\ndls_2}$, respectively. Then
$\cP^{[\deg\ndls_1+\deg\ndls_2]}$ is tamely ramified at
$\eta_{\ndls_1+\ndls_2}$.
\end{cor}

\begin{proof}
Put $d_i=\deg\ndls_i$ for $i=1,2$, and pull $\cP^{[d_i]}$ back to
$X_{\ndls_i}$ over the inverse image of $U^{(d_i)}$. Applying
\Cref{prop:ramification-external-product} to the two blowups
$X_{\ndls_i}\to C^{(d_i)}$, and writing $r_i$ for the morphisms supplied
by that proposition, shows that
\[
r_1^{-1}\cP^{[d_1]}\otimes^G r_2^{-1}\cP^{[d_2]}
\]
is tame; i.e., this contracted product has tame ramification at the generic
point of the exceptional divisor of
\[
\Bl_{(\ndls_1,\ndls_2)}
   \bigl(C^{(d_1)}\times_k C^{(d_2)}\bigr).
\]

By \Cref{lemma:takeuchi_lemma}, applied with
$\mdls=\ndls_1+\ndls_2$ and $d=d_1+d_2$, the addition morphism
\[
+_{d_1,d_2}\colon
C^{(d_1)}\times_k C^{(d_2)}
\longrightarrow C^{(d_1+d_2)}
\]
is \`etale at $(\ndls_1,\ndls_2)$ and maps this point to
$\ndls_1+\ndls_2$. Choose open neighborhoods
\[
(\ndls_1,\ndls_2)\in V\subset C^{(d_1)}\times_k C^{(d_2)},
\qquad
\ndls_1+\ndls_2\in W\subset C^{(d_1+d_2)}
\]
such that $+_{d_1,d_2}|_V\colon V\to W$ is \`etale and
\[
V\times_W\{\ndls_1+\ndls_2\}=\{(\ndls_1,\ndls_2)\}
\]
scheme-theoretically. Flat base change for blowups gives
\[
\Bl_{(\ndls_1,\ndls_2)}(V)
\cong
V\times_W\Bl_{\ndls_1+\ndls_2}(W).
\]
Hence the second projection defines an \`etale morphism
\[
\widetilde{+}_{d_1,d_2}\colon
\Bl_{(\ndls_1,\ndls_2)}(V)
\longrightarrow
\Bl_{\ndls_1+\ndls_2}(W)
\subset X_{\ndls_1+\ndls_2}.
\]
If $\eta$ denotes the generic point of the exceptional divisor of the
source, then
\[
\widetilde{+}_{d_1,d_2}(\eta)=\eta_{\ndls_1+\ndls_2}.
\]
By \Cref{prop:symmetric_powers_on_curves}, on the common dense open,
\[
(\widetilde{+}_{d_1,d_2})^{-1}\cP^{[d_1+d_2]}
\cong
r_1^{-1}\cP^{[d_1]}\otimes^G r_2^{-1}\cP^{[d_2]}.
\]
Its tameness therefore descends by
\Cref{lemma:descend_ramification_along_etale}.
\end{proof}

\noindent\textbf{From the point center to the degree-$d$ center.}

The comparison is governed by the addition morphism
\[
+\colon C^{(\deg\mdls)}\times_k C^{(d-\deg\mdls)}\longrightarrow C^{(d)}.
\]
After restricting the second factor to $U^{(d-\deg\mdls)}$ and blowing up,
it induces a comparison morphism
\[
X_{\mdls}\times_k U^{(d-\deg\mdls)}\longrightarrow X_{\mdls,d}.
\]
To descend tameness along this morphism, we must show that it is \`etale at
the generic point of the exceptional divisor in its source. This follows
from the case $l=1$ of the next lemma, adapted from
\cite[Lemma~4.1]{takeuchi2018blow}.

\begin{lemma}\label{lemma:takeuchi_lemma}
     Let $C$ be a projective, smooth, and geometrically connected curve over a perfect field $k$.
     Let $\mdls = \sum_{i=1}^{s} n_i P_i$ be an effective divisor where $P_1, \dots, P_s$ are distinct closed points.
     Let $U = C \setminus \mdls$ and let $d \geq \deg \mdls$.

     Suppose $\ndls_1, \dots, \ndls_l$ are pairwise coprime submoduli of $\mdls$ such that $\mdls = \sum_{j=1}^l \ndls_j$.
     Put $d_j=\deg\ndls_j$ for $1\leq j\leq l$ and $d_0=d-\deg\mdls$.
     Consider the summation morphism:
     \[
     \pi : Y:=C^{(d_1)} \times_k \dots \times_k C^{(d_l)} \times_k C^{(d_0)} \longrightarrow C^{(d)}
     \]
     defined by $(D_1, \dots, D_l, D_0) \mapsto \sum_{j=1}^l D_j + D_0$.

     Then $\pi$ is finite flat, and it is étale at the generic point of the closed subvariety
     \[
     V = \{\ndls_1\} \times_k \dots \times_k \{\ndls_l\} \times_k C^{(d_0)}\subset Y.
     \]
\end{lemma}

\begin{proof}
We may assume that $k$ is algebraically closed. The morphism $\pi$ is
finite: it is induced by the finite quotient morphism
\[
C^d/\bigl(S_{d_1}\times\cdots\times S_{d_l}\times S_{d_0}\bigr)
  \longrightarrow C^d/S_d.
\]
Its source and target are smooth $k$-schemes of pure dimension $d$.
Hence $\pi$ is flat by miracle flatness, and therefore finite flat.

Since the étale locus of $\pi$ is open and $V$ is integral, it is enough to
find one point of $V$ at which $\pi$ is étale. Choose an effective
divisor $D_0\in U^{(d_0)}$, and set
\[
v=\bigl(\ndls_1,\ldots,\ndls_l,D_0\bigr)\in V,
\qquad Q=\pi(v)=\mdls+D_0.
\]
The summands $\ndls_1,\ldots,\ndls_l,D_0$ have pairwise disjoint
supports, since the $\ndls_j$ are pairwise coprime and
$D_0$ is supported in $U$. Consider the quotient maps
\[
C^d\xrightarrow{\ p\ }
Y
\xrightarrow{\ \pi\ }C^{(d)},
\]
where $p$ groups the ordered points into blocks, and put $q=\pi\circ p$.
The same smoothness and miracle-flatness argument shows that $p$ and $q$
are finite flat, of degrees $h:=\prod_{j=0}^{l}d_j!$ and $d!$, respectively.
Write
\[
Q=\sum_{P\in\operatorname{Supp}(Q)}m_P P,
\qquad \mu:=\prod_{P\in\operatorname{Supp}(Q)}m_P!.
\]

The fiber $q^{-1}(Q)$ has $d!/\mu$ distinct points, corresponding to the
ordered lists representing $Q$. Its coordinate algebra $R$ has
$k$-dimension $d!$, since $q$ is finite flat of degree $d!$. Being
finite-dimensional over $k$, $R$ is Artinian. By \stackstag{00JB}, it
decomposes into its localizations at maximal ideals, which correspond to
the points of the fiber:
\[
R=\prod_{x\in q^{-1}(Q)}R_x,
\qquad R_x=\mathcal O_{q^{-1}(Q),x}.
\]
The action of $S_d$ is transitive on these points and induces isomorphisms
between the factors $R_x$. Therefore each has dimension
\[
\dim_k R_x=\frac{d!}{d!/\mu}=\mu.
\]

Let $A=\mathcal O_{\pi^{-1}(Q),v}$, so that $\Spec(A)$ is the local
part of the fiber at $v$, and let $B$ be the coordinate algebra of
$C^d\times_Y\Spec(A)$. As $p$ is finite flat of degree $h$, $B$ is a free
$A$-module of rank $h$. Hence
\[
\dim_k B=h\dim_k A.
\]
On the other hand, since the summands have disjoint supports, each point
$P$ of $Q$ occurs in exactly one block, with multiplicity $m_P$. The block
corresponding to $\ndls_j$ therefore has
$d_j!/\prod_{P\in\operatorname{Supp}(\ndls_j)}m_P!$ distinct orderings,
and the analogous formula holds for the block corresponding to $D_0$.
Multiplying these counts over all blocks gives exactly $h/\mu$ ordered
lists above $v$. Thus
\[
B\cong\prod_{\substack{x\in q^{-1}(Q)\\p(x)=v}}R_x,
\qquad
\dim_k B=\frac{h}{\mu}\,\mu=h.
\]
It follows that $\dim_k A=1$, and therefore $A=k$. The fiber of $\pi$ is
thus étale at $v$. Since $\pi$ is flat and of finite presentation, the
fiber criterion for étaleness (\stackstag{0476}) implies that $\pi$ is
étale at $v$, as required.
\end{proof}

\begin{cor}\label{cor:etale_at_generic_point}
For $d\geq\deg\mdls$, the addition morphism
\[
+\colon C^{(\deg\mdls)}\times_k C^{(d-\deg\mdls)}\longrightarrow C^{(d)}
\]
is finite flat, and it is \`etale at the generic point of
\[
\{\mdls\}\times_k C^{(d-\deg\mdls)}
   \subset C^{(\deg\mdls)}\times_k C^{(d-\deg\mdls)}.
\]
\end{cor}
\begin{proof}
This is the case $l=1$ and $\ndls_1=\mdls$ of \Cref{lemma:takeuchi_lemma}.
\end{proof}

We now pass from tameness on $X_{\mdls}$ to tameness on $X_{\mdls,d}$.

\begin{cor}\label{cor:tame-ramification-symmetric-product}
Let $\mdls$ be a modulus, let $d\geq\deg\mdls$, and let $\cP$ be a
$G$-torsor on $U$ whose ramification is bounded by $\mdls$. Assume that
$\cP^{[\deg\mdls]}$ is tamely ramified at the generic point
$\eta_{\mdls}$ of $E_{\mdls}$. Then $\cP^{[d]}$ is tamely ramified at
the generic point $\eta_0=\eta_{\mdls,d}$ of the exceptional divisor
$E_0=E_{\mdls,d}$ over $Z_0(\mdls,d)$.
\end{cor}
\begin{proof}
Work on $X_{\mdls}\times_k U^{(d-\deg\mdls)}$, with projections $p_1,p_2$,
and let $\theta$ be the generic point of
$E_{\mdls}\times_k U^{(d-\deg\mdls)}$. Stability of tame ramification
under base change (\stackstag{0EYD}) shows that
$p_1^{-1}\cP^{[\deg\mdls]}$ is tamely ramified at $\theta$. The torsor
$p_2^{-1}\cP^{[d-\deg\mdls]}$ is defined over the whole product, so it is
unramified, and hence tamely ramified, at $\theta$. By
\Cref{lemma:bounding_ramification_of_contracted_product},
\[
p_1^{-1}\cP^{[\deg\mdls]}
  \otimes^G p_2^{-1}\cP^{[d-\deg\mdls]}
\]
is therefore tamely ramified at $\theta$.

Consider the addition morphism
\[
+\colon C^{(\deg\mdls)}\times_k U^{(d-\deg\mdls)}\longrightarrow C^{(d)}.
\]
Its scheme-theoretic inverse image of $Z_0(\mdls,d)$ is
$\{\mdls\}\times_k U^{(d-\deg\mdls)}$. Flat base change for blowups
(\Cref{theorem:blowup_flat_base_change}) therefore gives the Cartesian
square
\[
\begin{tikzcd}[column sep=large,row sep=large]
X_{\mdls}\times_k U^{(d-\deg\mdls)}
  \arrow[r,"\widetilde{+}"] \arrow[d]
  & X_{\mdls,d} \arrow[d] \\
C^{(\deg\mdls)}\times_k U^{(d-\deg\mdls)}
  \arrow[r,"+"]
  & C^{(d)}.
\end{tikzcd}
\]
The vertical arrows are the blowup maps.
The induced morphism
$E_{\mdls}\times_k U^{(d-\deg\mdls)}\to E_0$ is flat, being a base
change of $+$. By going down for flat morphisms (\stackstag{00HS}), it
maps the generic point $\theta$ to the generic point $\eta_0$.
The left vertical map sends $\theta$ to the generic point of
$\{\mdls\}\times_k U^{(d-\deg\mdls)}$, where addition is étale by
\Cref{cor:etale_at_generic_point}. Since $\widetilde{+}$ is its base
change, it is étale at $\theta$.

On the dense open $U^{(\deg\mdls)}\times_k U^{(d-\deg\mdls)}$,
$\widetilde{+}$ agrees with addition. Hence
\Cref{prop:symmetric_powers_on_curves} gives
\[
(\widetilde{+})^{-1}\cP^{[d]}\cong
p_1^{-1}\cP^{[\deg\mdls]}
  \otimes^G p_2^{-1}\cP^{[d-\deg\mdls]}.
\]
This torsor is tamely ramified at $\theta$. Tameness therefore descends
along $\widetilde{+}$ to $\eta_0$ by
\Cref{lemma:descend_ramification_along_etale}.
\end{proof}


\begin{proof}[Proof of \Cref{theorem:SymmetricPowerOfSheafIsTamelyRamified}]
By flat base change, \Cref{lemma:weakly_unramified_base_change}, and
\Cref{prop:equivalence_ramification_bounded_by_modulus}, we may assume
that $k$ is algebraically closed.

For each single-point summand $\ndls=n_P P$ of $\mdls$,
\Cref{theorem:torsors_after_blowup_is_tame} gives tameness of
$\cP^{[\deg\ndls]}$ at $\eta_\ndls$. Applying
\Cref{cor:moduli-reduction} repeatedly gives tameness of
$\cP^{[\deg\mdls]}$ at $\eta_\mdls$.
\Cref{cor:tame-ramification-symmetric-product} then gives tameness of
$\cP^{[d]}$ at $\eta_0$ for every sufficiently large $d$. Since $H$ is a
nonempty open subset of $E_0$, it has generic point $\eta_0$, proving
the theorem over $\bar{k}$. The initial base-change reduction descends this
tameness to the original field $k$.
\end{proof}



% !TeX root = ../../../main.tex
\section[Extension of Tame Torsors to the Compactified Symmetric Power]{\texorpdfstring{Extension of Tame Torsors\\to the Compactified Symmetric Power}{Extension of Tame Torsors to the Compactified Symmetric Power}}\label{section:extension-across-boundary}

We now show that the tameness established in
\Cref{theorem:SymmetricPowerOfSheafIsTamelyRamified} implies extension
across $H$. Assume $\mdls\ne0$.

\begin{lemma}\label{lemma:ExtensionOfTamelyRamifiedSheaf}
Let $G$ be a finite abelian group, and let $\cP^{[d]}$ be a
$G$-torsor on $U^{(d)}$. If $\cP^{[d]}$ is tamely ramified at the generic
point $\eta_0$ of $H=\Cmod{d}\setminus U^{(d)}$, then it extends uniquely
to a $G$-torsor $\widetilde{\cP}^{[d]}$ on $\Cmod{d}$.
\end{lemma}

\begin{proof}
The scheme $\Cmod{d}$ is regular, being
a projective-space bundle over $\PicCm[d]$, which is smooth as a torsor under
$\PicCm[0]$. Its boundary $H$ is irreducible with generic point $\eta_0$ (see
the proof of \Cref{theorem:SymmetricPowerOfSheafIsTamelyRamified}). By purity
of the branch locus, it therefore suffices to prove that $\cP^{[d]}$ is
unramified at $\eta_0$.\footnote{The normalization of
$\Cmod{d}$ in $\cP^{[d]}$ is finite, since $\Cmod{d}$ is normal and
Noetherian and the generic cover is finite separable (\stackstag{032L}).
Unramifiedness at $\eta_0$ makes it \'etale in
codimension one, hence \'etale by purity (\stackstag{0BMB}). Full faithfulness
of restriction to $U^{(d)}$ (\stackstag{0BQI}) extends the $G$-action and
the torsor isomorphism uniquely, giving the unique $G$-torsor extension.}

Let $\xi$ be the generic point of $\PicCm[d]$, put
$\Omega=\kappa(\xi)^{\mathrm{sep}}$, and write $\bar\xi=\Spec\Omega$.
The relative hyperplane $H$ dominates $\PicCm[d]$, so
$\widetilde\Phi_d(\eta_0)=\xi$.
Thus $\eta_0$ is also the generic point of $H_\xi$. Since $\PicCm[d]$ is integral,
$\Oo_{\PicCm[d],\xi}=\kappa(\xi)$, so passing to the generic fiber does not
change its local ring:
\[
    \Oo_{\Cmod{d},\eta_0}=\Oo_{(\Cmod{d})_\xi,\eta_0}.
\]
Since $H_\xi$ is geometrically integral, $H_{\bar\xi}$ has a unique
generic point $\bar\eta$, which maps to
$\eta_0$ under the projection
$q:(\Cmod{d})_{\bar\xi}\to\Cmod{d}$.
Define the pulled-back $G$-torsor
\[
    \cP_\Omega:=\cP^{[d]}\times_{U^{(d)}}U^{(d)}_{\bar\xi}.
\]
\begin{samepage}
Every square in the following diagram is Cartesian:
\[
\begin{tikzcd}[column sep=huge, row sep=2em]
    \cP_\Omega \arrow[r] \arrow[d, purple]
      & \cP^{[d]} \arrow[d, purple] \\
    U^{(d)}_{\bar\xi} \arrow[r] \arrow[d, hook]
      & U^{(d)} \arrow[d, hook] \\
    (\Cmod{d})_{\bar\xi}
      \arrow[r, "q\colon\bar\eta\mapsto\eta_0"]
      \arrow[d, "(\widetilde\Phi_d)_{\bar\xi}"']
      & \Cmod{d} \arrow[d, "\widetilde\Phi_d"] \\
    \bar\xi=\Spec\Omega \arrow[r, "\bar\xi\mapsto\xi"']
      & \PicCm[d]
\end{tikzcd}
\]
\end{samepage}
By \Cref{lemma:geometric-generic-fiber-takeuchi},
\[
    \bigl((\Cmod{d})_{\bar\xi},H_{\bar\xi}\bigr)
    \cong(\mathbb P^r_\Omega,\mathbb P^{r-1}_\Omega),
    \qquad U^{(d)}_{\bar\xi}\cong\mathbb A^r_\Omega,
\]
where $r=d-\deg\mdls-g+1$.

Set
\[
    A=\Oo_{\Cmod{d},\eta_0},\qquad
    B=\Oo_{(\Cmod{d})_{\bar\xi},\bar\eta},\qquad
    K=\Frac(A),\quad L=\Frac(B).
\]
The projection $q$ induces $A\to B$, and the Cartesian diagram gives
\[
    \cP_\Omega|_{\Spec L}
    \cong\bigl(\cP^{[d]}|_{\Spec K}\bigr)
    \times_{\Spec K}\Spec L.
\]
The map $A\to B$ is obtained by extending scalars along the separable
algebraic extension $\kappa(\xi)\subset\Omega$; hence it is weakly
unramified with separable algebraic residue-field extension.
Hence \Cref{lemma:weakly_unramified_base_change} shows that
$\cP_\Omega$ is tamely ramified at $\bar\eta$ and gives
\[
    \cP^{[d]}\text{ is unramified at }\eta_0
    \quad\Longleftrightarrow\quad
    \cP_\Omega\text{ is unramified at }\bar\eta.
\]
Let
\[
    \chi:\pi_1(\mathbb A^r_\Omega)\longrightarrow G
\]
be the character corresponding to $\cP_\Omega$.
Since $G$ is abelian, $\chi$ factors as
\[
    \pi_1(\mathbb A^r_\Omega)
    \twoheadrightarrow\pi_1(\mathbb A^r_\Omega)^{\mathrm{ab}}
    \xrightarrow{\,\bar\chi\,}G.
\]
In characteristic zero, this abelianized fundamental group is trivial
\cite[Example~4.14]{tendler2015geometricclassfieldtheory}. Thus $\chi=1$,
so $\cP_\Omega$ is trivial and hence unramified.

Suppose that $\operatorname{char}(k)=p>0$. Then
$\pi_1(\mathbb A^r_\Omega)^{\mathrm{ab}}$ is pro-$p$
\cite[Example~4.15]{tendler2015geometricclassfieldtheory}, so
$\Gamma:=\operatorname{im}(\chi)$ is a finite $p$-group.
By \Cref{theorem:equivalence_fundamental_covers} and the discussion
preceding it, the connected cover corresponding to $\ker\chi$ is Galois
with group $\Gamma$, and $\cP_\Omega$ is a disjoint union of
$[G:\Gamma]$ copies of this cover.
Hence every field factor of $\cP_\Omega|_{\Spec L}$ is a Galois extension
of $L$ of degree $|\Gamma|$.

For each field factor, every ramification index $e$ with respect to the
valuation defined by $B$ satisfies
\[
    e\mid |\Gamma|,\qquad (e,p)=1
    \qquad\Longrightarrow\qquad e=1,
\]
where the prime-to-$p$ condition follows from tameness. Tameness also gives
separable residue-field extensions, so $\cP_\Omega$ is unramified at
$\bar\eta$.
\end{proof}



% !TeX root = ../../../main.tex
\section{Descent and Geometric Class Field Theory}\label{section:gcft-descent}

The preceding section shows that $\cP^{[d]}$ extends uniquely from
$U^{(d)}$ to Takeuchi's compactification $\Cmod{d}$. We now descend this
extension to $\PicCm[d]$. The essential point is that the compactified
Abel--Jacobi morphism has projective-space fibers: their trivial \'etale
fundamental groups imply that every finite \'etale torsor on the
compactification comes uniquely from the base.

When $\mdls=0$, the usual Abel--Jacobi morphism already has
projective-space fibers, so the same argument applies directly. We conclude
the proof of \Cref{thm:GCFT_reduced} by translating the
resulting torsor back into a rank-one local system.

\begin{lemma}[Descent along projective-space fibers]
\label{lemma:projective-space-fiber-descent}
Let $f\colon X\to S$ be a proper smooth morphism of finite presentation
over a connected scheme $S$, with geometric fibers isomorphic to projective
spaces. Then pullback induces an equivalence between finite \'etale covers
of $S$ and $X$, and hence between $G$-torsors on these schemes for every
finite group $G$.
\end{lemma}
\begin{proof}
For a geometric point $\bar s$ of $S$ and compatible geometric base points,
the homotopy exact sequence (\stackstag{0C0J}) gives
\[
\pi_1^{\mathrm{et}}(X_{\bar s})\longrightarrow
\pi_1^{\mathrm{et}}(X)\longrightarrow
\pi_1^{\mathrm{et}}(S)\longrightarrow 1.
\]
Since projective spaces over separably closed fields have trivial \'etale
fundamental group, $f$ induces an isomorphism
\[
\pi_1^{\mathrm{et}}(X)\xrightarrow{\sim}\pi_1^{\mathrm{et}}(S).
\]
The equivalence follows from the classification of finite \'etale covers by
finite sets with continuous fundamental-group action, and it respects
$G$-actions and the torsor condition.
\end{proof}

\begin{theorem}[Torsor formulation]\label{thm:GCFT_torsors}
Let $\Lambda$ be a finite commutative ring whose cardinality is invertible
in $k$, set $G=\Lambda^\times$, and let $\cP$ be a $G$-torsor on $U$ whose
ramification is bounded by $\mdls$. For every sufficiently large integer
$d$, there exists a $G$-torsor $\cQ_d$ on $\PicCm[d]$, unique up to
isomorphism, such that
\[
\Phi_d^*\cQ_d\cong\cP^{[d]}.
\]
\end{theorem}

\begin{proof}
The scheme $\PicCm[d]$ is geometrically connected, being a torsor under
the geometrically connected group $\PicCm[0]$.
When $\mdls=0$ and $d$ is sufficiently large, the usual Abel--Jacobi morphism
\[
\Phi_d\colon C^{(d)}\longrightarrow\PicC[d]
\]
is proper and smooth with projective-space geometric fibers, even when $C$
has no $k$-rational point
\cite[Corollary~3.9]{takeuchi2018blow}. Thus
\Cref{lemma:projective-space-fiber-descent} gives the required descent of
$\cP^{[d]}$, including uniqueness.

Suppose now that $\mdls>0$. For $d$ sufficiently large,
\Cref{theorem:takeuchi-compactification-theorem} extends $\Phi_d$ to a
projective space bundle
\[
\widetilde\Phi_d\colon
\Cmod{d}\longrightarrow\PicCm[d].
\]
By \Cref{theorem:SymmetricPowerOfSheafIsTamelyRamified},
$\cP^{[d]}$ is tamely ramified along
$H=\Cmod{d}\setminus U^{(d)}$. Hence
\Cref{lemma:ExtensionOfTamelyRamifiedSheaf} gives a unique extension
$\widetilde{\cP}^{[d]}$ to $\Cmod{d}$.

Applying \Cref{lemma:projective-space-fiber-descent} to
$\widetilde\Phi_d$, we obtain a $G$-torsor $\cQ_d$ on $\PicCm[d]$ with
\[
\widetilde\Phi_d^*\cQ_d\cong\widetilde{\cP}^{[d]}.
\]
Restriction to $U^{(d)}$ gives $\Phi_d^*\cQ_d\cong\cP^{[d]}$.
Full faithfulness of restriction from the normal scheme $\Cmod{d}$ to its
dense open $U^{(d)}$ (\stackstag{0BQI}), together with the descent
equivalence, shows that $\cQ_d$ is unique up to isomorphism with this
property.
\end{proof}

\begin{proof}[Proof of \Cref{thm:GCFT_reduced}]
Apply \Cref{thm:GCFT_torsors} to the frame torsor
$\cP=\operatorname{Fr}(\cF)$ under $G=\Lambda^\times$, and set
\[
\cG_d=\cQ_d\times^G\Lambda.
\]
The frame-torsor equivalence of \Cref{prop:torsor_module_equivalence} is
compatible with pullback and symmetric powers, so
$\Phi_d^*\cG_d\cong\cF^{[d]}$. The same equivalence transfers uniqueness
of $\cQ_d$ to uniqueness of $\cG_d$.
\end{proof}

